\section{Explicit logical macros in the reference language}
\label{app:old-macros}

This appendix fixes the finite macros used by the interpretations. All
operators in this appendix belong to the old conditional language
$\mathcal N$. They must be expanded there before applying the translation
$T$. Applying these definitions unchanged to arbitrary new mixed
closures would be incorrect. These definitions reproduce the interface
of \texttt{NativeGate.lean} and \texttt{NativeLogic.lean}
\cite{native}.

Write $C(v,a,b)=I(v,a;u.b)$, where $u$ is unused, and $1=S0$.
Bindings below are freshened so that no displayed passive parameter is
captured. Define an equation--formula selector $U_v(s,t;B)$ by
\[
\begin{aligned}
U_v(s,t;r=w)&=(C(v,s,r)=C(v,t,w)),\\
U_v(s,t;Q_y B)&=Q_y U_v(C(y,0,s),C(y,0,t);B),\\
U_v(s,t;D_{x,y} B)&=D_{x,y}U_v(C(y,0,s),C(y,0,t);B).
\end{aligned}
\]
Its two endpoints are the equation $s=t$ and the formula $B$.
Put $G_v A=U_v(0,0;A)$. Then $G_0 A$ is provable and
$G_{Sz}A\eqv A$. A selector between arbitrary formulas is
\[
\begin{aligned}
M_v(s=t,B)&=U_v(s,t;B),\\
M_v(Q_y A,B)&=Q_y M_v(A,G_yB),\\
M_v(D_{x,y} A,B)&=D_{x,y}M_v(A,G_yB).
\end{aligned}
\]
This recursion terminates on the depth of its first argument; the
separate recursion defining $U$ terminates on its formula argument.
For every $A,B$, the old calculus derives
$M_0(A,B)\eqv A$ and $M_{Sz}(A,B)\eqv B$.

Let $p(y)=I(y,0;u.C(u,1,0))$. The logical operations are
\[
\begin{aligned}
\bot_*&=(0=1),\\
A\Rightarrow_* B&=Q_y M_y(A,B),\\
A\wedge_* B&=Q_y G_y M_{p(y)}(A,B),\\
\exists_* x\,A(x)&=D_{x,y}G_y A(x),\\
A\vee_*B&=\exists_*v\,M_v(A,B),\\
\forall_*x\,A(x)&=A(0)\wedge_* Q_x A(x).
\end{aligned}
\]
All new bound variables are fresh. These are definitions by finite
syntax, not additional primitive logical axioms or rules.

\paragraph{Why they have the logical rules.}
Implication introduction uses the zero selector as its local hypothesis
and the successor selector as its conclusion; elimination specializes
at $1$. Existential introduction chooses the witness and uses the
successor gate. Existential elimination uses $D$ elimination, the
zero-gate theorem, and specialization at $1$; its final conclusion can
be any formula, independently of the existential matrix.

For conjunction, the gate supplies the base theorem and the computations
$p(1)=1$, $p(2)=0$ recover the two conjuncts. Numerical cases establish
the selector at an arbitrary control from proofs of both alternatives.
Disjunction uses witnesses $0$ and $1$; its elimination uses the derived
case rule on an implication to the desired conclusion. The universal
macro has its zero instance as one conjunct and every instance as the
other conjunct's closure. Finally, $0=1$ and the selector computations
give $s=C(0,s,t)=C(1,s,t)=t$ for arbitrary terms; structural induction
on formulas gives falsity elimination.

The full proofs and their substitution lemmas are preserved in the
pinned source files included with the package. The new theory's
normalization proof uses the translated versions of these derived-rule
constructions on $T$-image principal formulas before it defines the
ordinary logical operations on arbitrary new formulas.
